% !TEX program = xelatex
\documentclass[12pt]{ctexart}
\usepackage[a4paper,margin=2.2cm]{geometry}
\usepackage{enumitem}
\usepackage{amsmath,amssymb}
\setlist[itemize]{leftmargin=1.8em}

\title{凌泰炜硕士学位论文技术审查报告\\（重点：技术正确性、可复核性与公式严谨性）}
\author{paper-review skill}
\date{2026-04-08}

\begin{document}
\maketitle

\section*{总体评估}
论文围绕“基于头戴式毫米波雷达的第一人称人体运动姿态估计”展开，选题具有明确应用价值，章节组织总体完整，既提出了上肢动作估计模型和全身姿态估计模型，也实现了系统平台与实验验证。从工作量与论文结构看，已经具备硕士论文主体基础。

本次审查重点核查了第3--4章公式与实现描述、附录A中的实现型推导、参考文献与图表标签一致性。结论是：\textbf{论文当前的主要风险不在“有没有结果”，而在“公式是否写对、读者能否按文中定义直接复现实现”}。若仅从技术正确性与可复核性角度判断，当前稿件仍建议\textbf{大修（Major Revision）后再送审/归档}。

\section*{主要技术问题（按优先级）}

\subsection*{1) 6D旋转到旋转矩阵的 Gram--Schmidt 公式写法不保持 \(SO(3)\)}
\textbf{位置：} 第3章式(3.23)，第34页。\\
\textbf{问题：} 公式将旋转矩阵写成由三列向量组成，其中前两列分别写为归一化后的 \(\hat a\) 与
\(
\frac{b-(a\cdot b)\hat a}{\lVert b-(a\cdot b)\hat a\rVert}
\)，但第三列写成了 \(a\times b\)。若第三列仍使用未归一化的原始 \(a,b\)，则所得矩阵一般不再严格正交，也不保证列向量单位长度。\\
\textbf{影响：} 该式直接服务于关节旋转矩阵构造；若照当前公式实现，会削弱后续正运动学链条的几何一致性，属于会被答辩或送审专家直接追问的硬性问题。\\
\textbf{建议修复：}
\begin{itemize}
  \item 将第三列明确改为 \(\hat a\times \hat b\)（或对叉乘结果再次归一化）。
  \item 在式后补出完整定义：\(\hat a=a/\lVert a\rVert\)，\(\hat b=\frac{b-(a\cdot \hat a)\hat a}{\lVert b-(a\cdot \hat a)\hat a\rVert}\)。
  \item 明确说明最终矩阵满足 \(R_i^\top R_i=I\) 且 \(\det(R_i)=1\)。
\end{itemize}

\subsection*{2) 两处卷积公式按当前写法并未体现卷积核位移}
\textbf{位置：} 第3章式(3.16)，第32页；第4章式(4.10)，第51页。\\
\textbf{问题：} 两式都写成类似
\[
Q^{(l)}[m,n]=\sum_{k=0}^{K-1}W^{(l)}[m,k]\cdot Q^{(l-1)}[k,n-S\cdot P]+b^{(l)}[m]
\]
的形式，即求和变量 \(k\) 只出现在权重项中，没有进入输入索引的位移项。按该写法，每个 kernel tap 读取的是同一输入位置，严格说并不能构成可执行的卷积/空洞卷积。第4章对应公式存在同类问题。\\
\textbf{影响：} 这不是单纯记号“不好看”，而是会直接影响“他人能否按公式复现实现”的关键问题。\\
\textbf{建议修复：}
\begin{itemize}
  \item 明确输入索引应如何随 \(k\) 变化，例如写成 \(n+kP\)、\(n-S+kP\) 或与代码一致的等价形式。
  \item 明确定义 \(m,n,k\) 分别对应输出通道、空间位置、核内偏移还是其他轴，避免把通道维和空间维混写。
  \item 同时补充 padding、stride、dilation 的约定。
\end{itemize}

\subsection*{3) 第3章 BiLSTM 隐状态维度前后不一致}
\textbf{位置：} 第3章3.4.2，式(3.4)--(3.6)，第29页。\\
\textbf{问题：} 正文写明“隐藏层维度设为32，前向与反向传播分别生成隐藏状态，拼接后得到64维特征 \(h_{\mathrm{bilstm}}\in\mathbb{R}^{T_k\times T_s\times 64}\)”，但式(3.6)又把单向隐状态 \(\overrightarrow h_t,\overleftarrow h_t\) 直接写成 \(\mathbb{R}^{T_k\times T_s\times 64}\)。若每个方向已经是64维，则拼接后应为128维。\\
\textbf{影响：} 后续注意力层输入维数、参数矩阵大小和实现接口都会随之改变，属于明显的公式--实现接口不一致。\\
\textbf{建议修复：} 统一单向与双向输出维度的定义，建议明确写成“单向32维，拼接后64维”，并顺带检查后续 Attention/FC 的输入维度是否与之匹配。

\subsection*{4) 时间注意力汇聚公式缺少显式索引与归一化轴定义}
\textbf{位置：} 第3章式(3.7)，第29页；第4章式(4.18)，第53页。\\
\textbf{问题：} 两式都写成“\(\sum_i \mathrm{Softmax}(F(\cdot))*h\)”的形式，但被加权的特征项没有显式写成 \(h_i\)，Softmax 是沿时间维、特征维还是其他维度归一化也未说明。\\
\textbf{影响：} 这类注意力实现对归一化轴极其敏感。若只给出当前写法，审稿人无法判断这里到底是 temporal attention、self-attention，还是简单的加权池化。\\
\textbf{建议修复：}
\begin{itemize}
  \item 改写为分步形式，例如 \(\alpha_i=\mathrm{softmax}(f(h_i))\)，\(h_{\mathrm{atten}}=\sum_i \alpha_i h_i\)。
  \item 明确 \(\alpha_i\) 的归一化范围、形状以及输出维度。
\end{itemize}

\subsection*{5) 欧拉角转换公式存在象限歧义}
\textbf{位置：} 附录A式(A.16)，第78页。\\
\textbf{问题：} 偏航角与横滚角均写为 \(\arctan(\cdot/\cdot)\) 形式，而不是 \(\mathrm{atan2}(y,x)\)。\\
\textbf{影响：} 当分母接近0或姿态跨象限时，角度会发生 \(\pi\) 级跳变；这会影响姿态曲线展示、误差统计与复现实验的一致性。\\
\textbf{建议修复：}
\begin{itemize}
  \item 将相关反正切统一改为 \(\mathrm{atan2}\) 形式。
  \item 对 \(\arcsin\) 的输入补充截断说明（如 clamp 到 \([-1,1]\)），避免数值漂移导致非法输入。
\end{itemize}

\subsection*{6) IMU 测量模型与积分更新中的坐标变换方向未讲清}
\textbf{位置：} 附录A式(A.13)--(A.15)，第77--78页。\\
\textbf{问题：} 文中把 \(R_{WB}\) 定义为“IMU传感器到世界坐标系”的变换矩阵，但式(A.13)又用 \(R_{WB}(a_T-g)\) 去生成传感器测量加速度，式(A.14) 再用同一个 \(R_{WB}\) 把测量量积分到世界系。若 \(a_T\) 所在坐标系不另作声明，这里的方向约定是含混的。\\
\textbf{影响：} 这是附录中最容易被追问的“符号方向/坐标系约定”问题，直接关系到推导是否自洽。\\
\textbf{建议修复：}
\begin{itemize}
  \item 明确区分 \(R_{WB}\) 与 \(R_{BW}\)，并逐式标注向量所在坐标系。
  \item 明确 \(a_T\) 是世界系真值、机体系真值，还是去重力后的比力。
  \item 若正文实现采用标准 strapdown IMU 模型，建议按标准记号重写该段。
\end{itemize}

\subsection*{7) 实验可复核元信息仍然不足}
\textbf{位置：} 第3章3.5与第4章4.5，尤其表3.6--表3.10、表4.3--表4.4及图4-9/图4-10。\\
\textbf{问题：} 当前结果基本均为单次点估计；在已展示的数据采集与性能分析页中，尚未看到清晰的训练/验证/测试划分、跨被试隔离策略、重复运行次数、随机种子、优化器与关键超参数说明。\\
\textbf{影响：} 对深度学习模型而言，这会削弱结论的稳健性，也使结果难以被第三方复现。\\
\textbf{建议修复：}
\begin{itemize}
  \item 至少补充数据划分协议（按被试/按场景/按动作）与关键训练参数。
  \item 对核心表格补充多次运行的 mean\(\pm\)std，或至少说明是否固定随机种子并复验过稳定性。
\end{itemize}

\section*{有意义的编辑与规范问题}

\subsection*{1) 参考文献中存在完全重复条目}
\textbf{位置：} 参考文献 [8]/[97]、[24]/[55]、[73]/[77]。\\
\textbf{问题：} 以上三组条目在作者、标题、年份与出处上均为重复。\\
\textbf{影响：} 这会造成文中同一工作可能对应两个编号，削弱参考文献的规范性，也会让审稿人对 bibliography 管理过程产生质疑。\\
\textbf{建议修复：} 去重并统一正文引用编号；同时再整体扫一遍 bib 数据库，避免相同论文因大小写或缩写差异重复入库。

\subsection*{2) 基线模型名称存在重复拼写错误}
\textbf{位置：} 表3.6、表3.7、表3.8、表3.9、表4.3及相关图例。\\
\textbf{问题：} 多处将模型名写为“P4Tranformer”。若该基线对应 Transformer 模型，则此处应为“P4Transformer”。\\
\textbf{影响：} 模型名拼写错误会给人“表格是后期拼接而未统一校对”的印象，也不利于读者检索对应工作。\\
\textbf{建议修复：} 全文统一检查该模型名称，包括正文、表格、图例与附录说明。

\section*{快速校对扫尾（高置信、低成本）}
\begin{itemize}
  \item \textbf{图3-6，第28页：} 图内模块标签写成了“Attnetion”，应改为“Attention”。
  \item \textbf{图表图例多处：} “P4Tranformer”建议统一改为“P4Transformer”。
  \item \textbf{参考文献区：} 去重后重新顺排编号，并复核正文交叉引用是否全部更新到位。
\end{itemize}

\section*{建议优先修复顺序}
\begin{itemize}
  \item \textbf{第一优先级：} 先修正式(3.23)、式(3.16)、式(4.10)、式(A.16)、式(A.13)--(A.15) 等会直接影响技术正确性的公式。
  \item \textbf{第二优先级：} 统一第3--4章所有张量维度与注意力公式写法，使“公式、文字、实现接口”三者一致。
  \item \textbf{第三优先级：} 补足实验复现实用信息，至少说明数据划分、随机性控制与关键训练超参数。
  \item \textbf{第四优先级：} 清理重复参考文献与图表拼写错误，提升最终送审稿的规范性。
\end{itemize}

\section*{需要外部核验的项目}
\begin{itemize}
  \item 若参考文献中的 arXiv 预印本（如 [88]--[90] 一类条目）已经有正式发表版本，定稿时建议优先改引正式版本。
  \item 参考文献 [105] 目前以汇编卷中的二次出处形式给出；是否符合学校规范、是否应改引原始首发文献，建议作者按原文逐条核验。
\end{itemize}

\section*{结论}
这篇论文的研究问题与系统实现本身是有价值的，但当前稿件最需要补强的不是“再加几张结果图”，而是把关键公式、符号方向、维度定义和复现实用信息真正写严谨。若上述问题逐项修正，论文的技术可信度会明显提升；若不修，送审时很可能在公式正确性与可复核性上被集中追问。

\end{document}